%==============================================================
\section{Mở đầu}
\label{sec:intro}
%==============================================================

\subsection{Bối cảnh}

Điện toán lượng tử (quantum computing) là mô hình tính toán dựa trên các nguyên lý của cơ học lượng
tử: \emph{superposition}, \emph{entanglement} và \emph{interference}. Thay vì bit cổ điển chỉ nhận
giá trị $0$ hoặc $1$, đơn vị thông tin của máy tính lượng tử là \textbf{qubit} -- một vector đơn vị
trong không gian Hilbert hai chiều. Một hệ $n$ qubit được mô tả bởi $2^n$ amplitude phức, và các
thuật toán lượng tử khai thác cấu trúc này để giải một số bài toán với số bước ít hơn đáng kể so với
mọi thuật toán cổ điển đã biết: Grover tìm kiếm trong $O(\sqrt{N})$ bước thay vì $O(N)$
\cite{grover1996}, Shor phân tích thừa số nguyên tố trong thời gian đa thức \cite{shor1997}.

Tuy nhiên, các máy tính lượng tử hiện nay thuộc thời kỳ \textbf{NISQ} (Noisy Intermediate-Scale
Quantum) \cite{preskill2018nisq}: vài chục tới vài trăm qubit, chưa có sửa lỗi đầy đủ, và mỗi phép
toán đều có sai số. Cụ thể, một thiết bị thật áp đặt ba ràng buộc lên mạch lượng tử:
\begin{enumerate}[noitemsep]
  \item \textbf{Basis gates} -- chỉ thực hiện được một tập cổng nhỏ, ví dụ
    $\{\code{rz}, \code{sx}, \code{x}, \code{cx}\}$.
  \item \textbf{Connectivity} -- cổng hai qubit chỉ thực hiện được trên các cặp qubit nối trực tiếp
    với nhau theo một đồ thị cố định (\emph{coupling map}).
  \item \textbf{Noise} -- cổng hai qubit có xác suất lỗi khoảng $10^{-3}$--$10^{-2}$, tệ hơn cổng
    một qubit một tới hai bậc.
\end{enumerate}

Người lập trình viết mạch ở mức trừu tượng (dùng Toffoli, SWAP, cổng giữa hai qubit bất kỳ), còn
thiết bị chỉ chạy được mạch thoả cả ba ràng buộc trên. Bước dịch mạch lý tưởng thành mạch chạy được
gọi là \textbf{quantum circuit compilation} -- trong Qiskit gọi là \textbf{transpilation}. Hình
\ref{fig:motivation} minh hoạ mức độ chênh lệch: một mạch 5 qubit có độ sâu (\emph{depth}) 6 sau khi
biên dịch lên một thiết bị có coupling map dạng đường thẳng trở thành mạch có depth 40 với 19 cổng
\cx{}. Toàn bộ khoảng cách đó do compiler quyết định, và mỗi cổng \cx{} tiết kiệm được là một phần
sai số bớt đi.

\begin{figure}[h]
  \centering
  \begin{tikzpicture}[node distance=1.2cm]
    \node[box=AccentBlue, minimum width=3.4cm, minimum height=1.6cm] (src)
      {\textbf{Mạch lý tưởng}\\ depth 6, size 11\\ \code{ccx}, \code{swap}, \code{cz}, \code{ry}};
    \node[box=Purple, right=1.4cm of src, minimum width=2.8cm, minimum height=1.6cm] (tr)
      {\textbf{Transpiler}\\ layout, routing,\\ synthesis, tối ưu};
    \node[box=ProbRed, right=1.4cm of tr, minimum width=3.4cm, minimum height=1.6cm] (dst)
      {\textbf{Mạch chạy được}\\ depth 40, size 65\\ 19 \code{cx}, 27 \code{rz}, 12 \code{sx}};
    \draw[arr] (src) -- (tr);
    \draw[arr] (tr) -- (dst);
    \node[below=0.25cm of tr, font=\footnotesize, align=center, text=black!70]
      {basis $\{\code{cx},\code{rz},\code{sx},\code{x}\}$, coupling map \code{line(5)}};
  \end{tikzpicture}
  \caption{Cùng một thuật toán, mạch thật sự chạy trên thiết bị lớn hơn mạch được viết rất nhiều
  (số liệu đo trong notebook \code{02\_compilation/00}, Qiskit 2.5.2, \code{optimization\_level=2}).}
  \label{fig:motivation}
\end{figure}

Vì phần lớn các bài toán con của compilation (chọn layout, chèn SWAP tối thiểu, tổng hợp cổng tối
ưu) đều là NP-hard \cite{siraichi2018allocation}, các compiler thực tế dùng \emph{heuristic} do con
người thiết kế, tiêu biểu là thuật toán SABRE \cite{li2019sabre}. Điều này mở ra một hướng nghiên cứu
đang phát triển mạnh: dùng \textbf{reinforcement learning} (RL) để học các quyết định biên dịch tốt
hơn heuristic \cite{pozzi2022rlrouting, fosel2021rlopt, ruiz2025alphatensor, riu2025zx}.

\subsection{Mục tiêu của chuyên đề}

Chuyên đề nghiên cứu 1 đặt mục tiêu xây dựng nền tảng đủ vững để bắt đầu nghiên cứu về biên dịch mạch
lượng tử, cụ thể:
\begin{enumerate}[noitemsep]
  \item Nắm vững các khái niệm cơ bản: qubit, ký hiệu Dirac (ket/bra), superposition, phép đo,
    entanglement.
  \item Hiểu các cổng lượng tử một và nhiều qubit dưới dạng ma trận unitary, và cách chúng tạo nên
    mạch lượng tử.
  \item Hiểu cơ chế chung của các thuật toán lượng tử cơ bản (phase kickback, interference) qua
    Deutsch--Jozsa, Bernstein--Vazirani, Grover, teleportation và superdense coding.
  \item Phân tích quá trình biên dịch mạch trong Qiskit như một \textbf{bài toán tối ưu}: các stage,
    cấu trúc dữ liệu, hàm cost, và các thuật toán layout, routing, synthesis.
  \item Phát biểu qubit routing thành một Markov decision process (MDP), hiện thực thành môi trường
    \code{gymnasium} \textbf{đúng semantics}, và dựng \textbf{baseline trung thực} để làm nền cho
    các thí nghiệm RL ở Chuyên đề nghiên cứu 2.
\end{enumerate}

\subsection{Phạm vi và công cụ}

Báo cáo giới hạn ở mô hình mạch lượng tử (quantum circuit model) và các thiết bị dạng cổng
(gate-based). Mọi thí nghiệm được thực hiện bằng \textbf{Qiskit 2.5} \cite{qiskit2024} và
\code{qiskit-aer} trên simulator, với \code{gymnasium} cho môi trường RL. Toàn bộ mã nguồn được tổ
chức thành một chuỗi Jupyter notebook trong kho mã \code{master-research-1}, chia làm ba phần tương
ứng với các chương của báo cáo (Phụ lục~B). Mọi con số trong báo cáo đều tái lập được
bằng notebook tương ứng.

\subsection{Cấu trúc báo cáo}

Báo cáo đi theo trình tự \emph{ket $\to$ cổng $\to$ thuật toán $\to$ compilation}, mỗi chương là nền
cho chương sau:
\begin{itemize}[noitemsep]
  \item \textbf{Mục~\ref{sec:basics}} trình bày qubit, ký hiệu Dirac, superposition, phép đo và hệ
    nhiều qubit.
  \item \textbf{Mục~\ref{sec:gates}} trình bày các cổng lượng tử, entanglement và quy trình làm việc
    với Qiskit.
  \item \textbf{Mục~\ref{sec:algorithms}} trình bày các thuật toán lượng tử cơ bản xoay quanh cơ chế
    phase kickback.
  \item \textbf{Mục~\ref{sec:compilation}} phân tích quá trình biên dịch: pipeline, layout, routing,
    synthesis và cách tự viết pass.
  \item \textbf{Mục~\ref{sec:mdp}} phát biểu compilation như một bài toán MDP cho reinforcement
    learning.
  \item \textbf{Mục~\ref{sec:experiments}} tổng hợp các kết quả thực nghiệm và bài học phương pháp.
  \item \textbf{Mục~\ref{sec:conclusion}} kết luận và đề xuất hướng phát triển.
\end{itemize}
